% This is LLNCS.DEM the demonstration file of
% the LaTeX macro package from Springer-Verlag
% for Lecture Notes in Computer Science,
% version 2.3 for LaTeX2e
%
\def\techreport{true}
\ifx\techreport\undefined
\documentclass{llncs}
\else
\documentclass[10pt]{article}
\fi

%
\usepackage{makeidx}  % allows for indexgeneration
\usepackage[pdftex]{graphicx}
\usepackage{multicol}
\usepackage{url}
\usepackage{amsmath}
\usepackage{amssymb}
\usepackage{stmaryrd}
\usepackage{algorithmicx}
\usepackage[ruled]{algorithm}
\usepackage{algpseudocode}
%\usepackage{latex8}
%

% Make floats work with multicols (maybe)
\makeatletter
\newenvironment{figurehere}
	{\def\@captype{figure}}
	{}
\makeatother


\begin{document}
%
%\frontmatter          % for the preliminaries
%
\ifx\techreport\undefined
\pagestyle{headings}  % switches on printing of running heads
\addtocmark{Incremental Cycle Analysis} % additional mark in the TOC
%
%\tableofcontents
%
\mainmatter              % start of the contributions
\fi
%
\title{Towards Incremental Cycle Analysis in ESMoL Distributed Control System
Models}
%
\ifx\techreport\undefined
\titlerunning{Incremental Cycle Analysis}  % abbreviated title (for running head)
%                                     also used for the TOC unless
%                                     \toctitle is used
%
\author{Joseph Porter \and Daniel Balasubramanian \and Graham Hemingway \and J\'{a}nos Sztipanovits}
%
\authorrunning{Joseph Porter et al.}   % abbreviated author list (for running head)
%
%%%% list of authors for the TOC (use if author list has to be modified)
\tocauthor{Joseph Porter}
%
\institute{Institute for Software Integrated Systems, \\Vanderbilt University, \\Nashville TN 37212, USA,\\
\email{jporter@isis.vanderbilt.edu},\\ WWW home page:
\texttt{http://www.isis.vanderbilt.edu}
}
\fi

\ifx\techreport\undefined
\else
\begin{figure}[!htp]
\label{fig:ISIS} \vspace{1cm}
\begin{center}
\includegraphics[width=3in,height=2in]{figures/isis_old.png}
\end{center}
\vspace{.0cm}
\end{figure}

\begin{center}\begin{Large}
Institute for Software Integrated Systems \\
Vanderbilt University \\
Nashville, Tennessee, 37235 \\
\end{Large}\end{center}

\vspace{1cm}

\begin{LARGE}
\begin{center}
Towards Incremental Cycle Analysis in ESMoL Distributed Control System Models \\
\vspace{.75cm}
\begin{Large}
Joseph Porter, Daniel Balasubramanian, Graham Hemingway, and J\'{a}nos Sztipanovits
\end{Large}\\
\end{center}
\end{LARGE}

\vspace{2cm}

\begin{Large}
\begin{center}
\textbf{TECHNICAL REPORT} \\
\ \\
ISIS-11-106
\end{center}
\end{Large}

\newpage

\fi

\ifx\techreport\undefined
\maketitle % typeset the title of the contribution
\fi


\begin{abstract}
We consider the problem of incremental cycle analysis for dataflow models in 
the Embedded Systems Modeling Language (ESMoL).  We give a general form of a
cycle enumeration algorithm that makes use of graph hierarchy to improve
analysis efficiency.  Our framework also stores simple connectivity information
in the model to accelerate future cycle analyses when additional components
are added or modifications are made.  Finally we give a mapping from a term
algebraic model of the ESMoL component model and logical dataflow sublanguages
to the analysis framework, and an evaluation on a fixed-wing aircraft controller model.  This is part of a larger effort to integrate cycle analysis into the ESMoL tool suite to aid well-formedness checking during model construction.
\ifx\techreport\undefined
\keywords{model-based design, embedded systems, model analysis}
\fi
\end{abstract}
%
%\twocolumn
%\begin{multicols}{2}


Embedded software often operates in environments critical to human life and subject to our direct expectations. We assume that a handheld MP3 player will perform reliably, or that the unseen aircraft control system aboard our flight will function safely and correctly. Safety-critical embedded environments require far more care than provided by the current best practices in software development. Embedded systems design challenges are well-documented~\cite{HenSif:2006}, but industrial practice still falls short of these expectations for many kinds of embedded systems.

In modern designs, graphical modeling and simulation tools (e.g. Mathworks' Simulink/Stateflow) represent physical systems and engineering designs using block diagram notations. Design work revolves around simulation and test cases, with code generated from "`complete"' designs. Control designs often ignore software design constraints and issues arising from embedded platform choices. At early stages of the design, platforms may be vaguely specified to engineers as sets of tradeoffs.

Software development uses UML (or similar) tools to capture concepts such as components, interactions, timing, fault handling, and deployment. Workflows focus on source code organization and management, followed by testing and debugging on target hardware. Physical and environmental constraints are not represented by the tools. At best such constraints may be provided as documentation to developers.

Complete systems rely on both aspects of a design.  
Designers lack tools to model the interactions between the
hardware, software, and the environment.  For example, software 
generated from a carefully simulated functional dataflow model may 
fail to perform correctly when its functions are distributed over a 
shared network of processing nodes.  Cost considerations may force the
selection of platform hardware that limits timing accuracy.  Neither 
aspect of development supports comprehensive validation of certification 
requirements to meet government safety standards. 

We propose a suite of tools
that aim to address many of these challenges.  Currently under development at
Vanderbilt's Institute for Software Integrated Systems (ISIS), these 
tools use the Embedded Systems Modeling Language (ESMoL), which is a
suite of domain-specific modeling languages (DSML) to integrate the 
disparate aspects of a safety-critical embedded systems design and 
maintain proper separation of concerns between engineering and software 
development teams.  Many of the concepts and features presented here also 
exist separately in other tools.  We describe a model-based approach to building 
a unified model-based design and integration tool suite which has the potential to go far beyond the state of the art.

In the sequel we will provide an overview of the tool vision, and then describe the 
features of these tools from the point of view of available functionality.   
Note that two different development processes will be discussed -- the development of a distributed control system implementation (by an imagined user of the tools), and our development of the tool suite itself.  The initial vision section illustrates how the tools would be used to model and develop a control system.  The final sections describe different parts of our tool-development process in decreasing order of maturity.  We strive for clarity, with an apology to the diligent reader where the distinction is unclear.


\section{Background}
\label{section:background}

\subsection{ESMoL Component Model}

As the ESMoL language structure is documented elsewhere\cite{modeling:esmol_tr}, we only cover details relevant to incremental cycle checking.  ESMoL is a graphical
modeling language which allows designers to use Simulink diagrams as 
synchronous software function specifications (where the execution of each 
block is equivalent to a single bounded-time blocking C language call).  These
specifications are used to create blocks representing ESMoL component types.   
ESMoL components have message structures as interfaces, and the type
specification includes a map between Simulink signal ports and the 
fields of the input and output message structures.  The messages represent C structures, and the map graphically captures the marshaling of Simulink data to those structures.

Once software component types and interfaces have been specified, ESMoL 
designers instantiate those components into a distributed 
deployment model.  ESMoL allows the separate specification of the logical
data flow, the mapping of component instances to hardware, timing
information for tasks executing those components, and timing for messages sent
over a time-triggered communication bus.  Code generated from the models
conforms to an API for time-triggered execution.  A portable virtual machine
implementation of the API allows execution in simulation, hardware-in-the-loop,
and fully deployed configurations\cite{modeling:truetime}.

ESMoL deliberately provides an unusual degree of freedom in creating 
software component types.  A designer can include Simulink references from
any part of an imported dataflow model, and instantiate them any number of 
times within the type definitions. The partition of functions into ESMoL types 
allows the designer to control the granularity of functions assigned to 
distributed tasks.  Tasks can distribute functions over a time-triggered 
network for performance, or replicate similar functions for fault mitigation.  
This level of flexibility requires automatic type-checking to ensure 
compatibility for chosen configurations.  Beyond interface type-checking, 
structural well-formedness problems arise during assembly such as zero-delay 
cycles.   Model analysis must ensure well-formedness.

\subsection{Cycle Enumeration}

To implement cycle enumeration we use the algorithm Johnson proposed as an 
extension of Tiernan's algorithm \cite{cycles:tiernan} for 
enumerating elementary cycles in a directed graph\cite{cycles:johnson75}.  Both
approaches rely on depth-first search with backtracking, but Johnson's method
marks vertices on elementary paths already considered to eliminate fruitless
searching, unmarking them only when a cycle is found.  Johnson's algorithm is 
polynomial ($O((n + e)c)$, where $n$, $e$, and $c$ are the sizes of the 
vertex, edge, and cycle set, respectively), and is still considered the best 
available general cycle enumeration method\cite{cycles:mateti76}.  We created
an implementation of Johnson's algorithm in C++ using the Boost Graph library
\cite{tools:boostgraph}.

\subsection{Hierarchical Graphs}

For formally describing our incremental approach we use the algebra of
hierarchical graphs introduced by Bruni et al\cite{graphs:hier_algebra}.
We repeat here their first definition: a \emph{design} is a term of
sort $\mathcal{D}$ generated by

\begin{align}
\label{eq:hiergraphs}
\mathbb{D} &::= L_{\bar{x}} [ \mathbb{G} ]  \\
\mathbb{G} &::= \mathbf{0} \mid x \mid l<\bar{x}> \mid \mathbb{G} \parallel \mathbb{G} \mid  (\nu \bar{x})\mathbb{G} \mid \mathbb{D} <\bar{x}>  \nonumber
\end{align}

Here term $\mathbb{G}$ represents a hierarchical directed graph, $\mathbb{D}$ 
is an edge-encapsulated hierarchical graph, $x$ is a vertex,
$\bar{x}$ is a list of vertices in $\mathbb{G}$ (for which 
$\lfloor \bar{x} \rfloor$ is the corresponding set), $l \in \mathcal{E}$ 
(edge labels of $\mathbb{G}$, where edges can have n-ary connectivity ), 
$L_{\bar{x}} \in \mathcal{D}$ ($\mathcal{D}$ are
the design labels of $\mathbb{G}$ and $\bar{x}$ 
are interface vertices in $L$),
$\mathbb{G} \parallel \mathbb{G}$ is parallel graph composition which merges 
vertices with common names, $(\nu \bar{x}) \mathbb{G}$ restricts the interface 
of graph $\mathbb{G}$ to exclude vertices in $\lfloor \bar{x} \rfloor$,
and the notation $\mathbb{D} <\bar{x}>$ maps the vertices from the 
interface of $\mathbb{D}$ to the vertices listed in $\bar{x}$ (renaming
vertices internal to the design for the external interface). Finally
$\llbracket \mathbb{G} \rrbracket$ indicates the graph corresponding
to the term $\mathbb{G}$.

Intuitively, Equation \ref{eq:hiergraphs} is a grammar defining a simple textual notation for describing typed hierarchical graphs.  Within the formalism we can compare equivalence between algebraic descriptions of two hierarchical graphs using reduction rules and a normal form (as in Bruni\cite{graphs:hier_algebra}), though equivalence is beyond the scope of this publication. The algebraic properties are for future use. The other main attraction of this particular formalism is that the notation allows the definition of interface symbols which correspond easily to port objects in a dataflow language, and the hiding of those interfaces as we specialize types.  The notation is a compact shorthand for much larger diagrams or mathematical descriptions. The design sorts $\mathbb{D}$ correspond to composite types in our dataflow language (which may have children), and the edge-encapsulation means that graph edges are not visible outside the component that contains them.   The specification for a composite element is $L_{\bar{x}} [\mathbb{G} ]$, which means that an element of the sort $\mathbb{D}$ has type $L$ with interface vertices in the list $\bar{x}$ and a corresponding internal graph $\mathbb{G}$ defining the details of the component. The internal graph may also include subcomponents. Gluing of subgraphs (contained by the design sorts $\mathbb{D}$) only occurs at common vertices.  When a composite element from $\mathbb{D}$ of a particular type is used as a child element to form a  larger (parent) graph, vertices from the child are possibly renamed in the parent, hence the notation $\mathbb{D} <\bar{x}> $.  In a parallel composition, vertices with the same name $x$ are glued together.

Note that the term \emph{design} in the formalism is synonymous with the concept of a \emph{component} in a modeling hierarchy, with a type and a set of interface vertices. Unfortunately in the realm of graph theory the term \emph{component} has a different meaning.  The term-algebraic graph formalism also includes a definition of \emph{well-typedness}, where types defined on the vertex set are only connected if their types are compatible.  Finally the authors define \emph{well-formedness} for hierarchical graphs which includes \emph{well-typedness} as a condition.  We do not define the entire formalism here, only enough to understand the essence of the 
connections between the terms and the graphs that they represent.



\section{Incremental Cycle Analysis}

Our intention is to support a design and analysis work flow that includes 
incremental analysis steps.  For example, a designer may analyze part of the design before integrating it into a larger part of the system.  In our work flow, we envision storing the results of that first analysis along with some interface data to reduce the cost of the second analysis.  The same should hold true for the system design.  We should be able to analyze the system design efficiently, calculating incremental analysis interfaces.  When the system models are revised, whether by adding, removing, or modifying components we can isolate the effects of the change on the cost of the analysis.  Cycle analysis is a useful example, but our aim is to tackle this problem more generally.

\subsection*{Formal Model}
\addcontentsline{toc}{subsection}{Formal Model}

Let $\mathbb{G}$ be a well-formed hierarchical graph (as in Bruni
\cite{graphs:hier_algebra}).  To get more comfortable with the notation,
first note that graph $\mathbb{G}$ itself (without hierarchical structure)
can be given as:

\begin{equation}
\mathbb{G} = (\parallel x) \parallel (\parallel_{(u,v) \in \mathcal{E}} l < u, v >)
\end{equation}

which is the parallel composition of the graphs induced by individual edges of $\mathbb{G}$, merged at their common vertices.

Let $C(\mathbb{G})$ be the set of elementary cycles in $\mathbb{G}$, and 
let $P(\mathbb{G}, u, v)$ be the subgraph of $\mathbb{G}$ containing all of
the paths from vertex $u$ to vertex $v$.

Consider a design of type $W$.  Let $W_{\bar{x}}^p [G]$ 
represent a parent design object in a graph hierarchy with interface
vertices $\bar{x}$, 
and let $W_{\bar{x}_i}^{c_i} [G]$ be the children
of $W^p$ 
($\llbracket W^{c_i} \rrbracket \subset \llbracket W^p \rrbracket$).  
Then neglecting vertex hiding and renaming 
to simplify the illustration, we have the following:

\begin{equation}
W_{\bar{x}}^p [\mathbb{G}] = W_{\bar{x}}^p [ (\parallel_{h} x_h) \parallel
(\parallel_{(j,k)} l< j,k >) \parallel (\parallel_{i} W_{\bar{x}_i}^{c_i} 
[\mathbb{G}]) ]
\label{eq:composition}
\end{equation}

Eq. \ref{eq:composition} describes the design $W^p$ in terms of its design 
children $W^{c_i}$, internal vertices $x_h$, and edges $l <j,k>$.   Fig. \ref{fig:hier_cycle} gives a simple example of a hierarchical graph with one child component and one cycle including vertices and edges within the child component.

\begin{figure}[htb]	
    \begin{center}
    \centerline{\includegraphics[width=0.7\columnwidth]{figures/original_graph.png}}
    \caption{Generic hierarchical graph model}
    \label{fig:hier_cycle}
    \end{center}	
\end{figure}

We introduce a new label $l_c$ into the sort for edges ($\mathcal{E}$),  
which is used to connect vertices at the boundaries of a design, 
abstracting the interface connectivity of the design. Introduce 
a new mapping $A: \mathcal{D} \rightarrow \mathcal{D}'$ from the 
designs of $\mathbb{G}$ to designs in a new graph $\mathbb{G}'$.  
$\mathbb{G}'$ is identical to $\mathbb{G}$, but adds the new edge
label.  This is the interface that we will use for incremental cycle
analysis.

\begin{align}
A( W_{\bar{x}_i}^{c_i} [\mathbb{G}] ) = & W_{\bar{x}_i}^{c_i} 
[ (\parallel_h x_h) \parallel 
(\parallel_{(j,k) \in \lfloor \bar{x}_i \rfloor \wedge P(\mathbb{G},j,k) 
\neq \emptyset} l_c <j,k>) \\
 & \parallel (\parallel_{(j,k)} l<j,k> )
\parallel (\parallel_m W_{\bar{x}_m}^{c_m} [\mathbb{G}]) ]) ] \nonumber
\end{align}

In this abstraction function the child designs are replaced by a much simpler
connectivity graph.  We introduce two functions to support the algorithm:

\begin{align}
&R(A(W_{\bar{x}}^{c_i} [\mathbb{G}])) = W_{\bar{x}_i}^{c_i} [ (\parallel_{x \in \bar{x}_i} x) \parallel (\parallel_{(j,k) \in l_c} l_c<j,k> ) ] \\
&S(W_{\bar{x}}^p [\mathbb{G}]) = W_{\bar{x}}^p [ (\parallel_h x_h) \parallel
(\parallel_{(j,k) \in \lfloor \bar{x} \rfloor} l<j,k>) \parallel
(\parallel_{i} R(A(W_{\bar{x}_i}^{c_i} [\mathbb{G}]))) ] 
\end{align}

$R(\cdot)$ and $S(\cdot)$ map designs in $\mathbb{G}$ to 
an abstracted design which only has connectivity edges 
for each child design.  In other words, when analyzing a component
of $\mathbb{G}$ we use the incremental interface data for each child
component rather than its full details. This is a useful abstraction for 
cycle detection: we can exploit the graph hierarchy to 
enumerate simple cycles more efficiently.  Figure \ref{fig:abs_cycle} gives an example of the transformation defined by $A(\cdot)$, $R(\cdot)$, and $S(\cdot)$.  The child graph is replaced by its abstracted correspondent, which only preserves connectivity between interface vertices.

\begin{figure}[htb]	
    \begin{center}
    \centerline{\includegraphics[width=0.7\columnwidth]{figures/abstracted_graph.png}}
    \caption{Abstract hierarchical graph model}
    \label{fig:abs_cycle}
    \end{center}	
\end{figure}

\subsection*{Algorithm Description}
\addcontentsline{toc}{subsection}{Algorithm Description}

Assume we have a function 
${\scriptstyle \mathrm{FINDALLCYCLES}}: \mathbb{G} \rightarrow \mathbf{2}^{\mathbb{G}}$
which enumerates all elementary cycles in a graph $\mathbb{G}$, 
returning sets of subgraphs.  Then Algorithm \ref{alg:hcycles} 
adapts the general algorithm
${\scriptstyle \mathrm{FINDALLCYCLES}}$ to the hierarchical graph 
structure described above.  We 
assume that $\mathbb{G}$ has a unique root design, and that we have a 
function $modified: \mathbb{D} \rightarrow boolean$ which indicates whether
a particular hierarchical component has been modified since the last run.
New components in the model are considered modified by default.

\begin{algorithm}[H]
\caption{Hierarchical cycle detection}
\label{alg:hcycles}
\begin{algorithmic}[1]
\State $cycles \gets [ ]$
\State $ifaces \gets \{ \}$
\Function{findhcycles}{ $ \llbracket W_{\bar{x}}^p [\mathbb{G}] \rrbracket$ }
   \ForAll { $W_{\bar{x}_i}^{c_i} [\mathbb{G}] \in W_{\bar{x}}^p [\mathbb{G}]$  }
      \State ${\scriptstyle \mathrm{FINDHCYCLES}}( \llbracket W_{\bar{x}_i}^{c_i} [\mathbb{G}] \rrbracket )$
   \EndFor
   \State $modified( W_{\bar{x}}^p [\mathbb{G}]) \gets (modified( W_{\bar{x}}^p [\mathbb{G}]) \vee ( \vee_{c_i} modified( W_{\bar{x}_i}^{c_i} [\mathbb{G}] ) )$
   \If { $modified( W_{\bar{x}}^p [\mathbb{G}])$ }
      \State $T \gets \llbracket S(W_{\bar{x}}^p [\mathbb{G}] \rrbracket )$
      \State $cycles \gets [ cycles; {\scriptstyle \mathrm{FINDALLCYCLES}}( T ) ]$
      \State $ifaces[p] \gets A(T)$
   \EndIf
\EndFunction
\State ${\scriptstyle \mathrm{FINDHCYCLES}}( \mathbb{G}) $
\end{algorithmic}
\end{algorithm}


\begin{itemize}
\item (Lines 1-2) We initialize a list to contain the resulting cycles, and an associative list to contain the interface data.
\item (Lines 3-6) The algorithm performs a depth-first search on the hierarchical graph, recursively visiting all of the child components ($W_{\bar{x}_i}^{c_i} [\mathbb{G}]$)  for the current (parent) component ($W_{\bar{x}}^p [\mathbb{G}]$) which is currently under evaluation.
\item (Line 7)  The modification status is propagated up the hierarchy as the algorithm progresses.  Each component which has a modified child will also be marked as modified.
\item (Lines 8-12) If the current component has been modified, we use the previously computed incremental connectivity interface for each subcomponent to check for cycles in the parent component -- the connectivity graph interface is substituted for each subcomponent.  The cycles are accumulated as the algorithm ascends to the top of the model, and a connectivity interface is created for the current component as well.
\end{itemize}

The runtime for the extended algorithm should be slightly worse than Johnson's 
algorithm in the worst case, as it must also compute the interface graphs.  
In the average case the cycle checking proceeds on graphs much smaller than
the global graph, offsetting the cost of finding paths in each subgraph.
Further, if the incremental interface edges are stored in the model following
the analysis, then scalability is enhanced when incrementally adding functions to a design.  Cycle analysis is then restricted to the size of the new components together with the stored interfaces.
 

\section{ ESMoL Language Mapping }
\label{section:esmol}

%Now to map ESMoL logical architecture models onto this cycle-checking
%formal model we use the following rules:
%
%\begin{align}
%\mathbf{Subsys} ::=\; & L_{\bar{i},\bar{o}}^{Subsys} \llbracket \mathbf{Dataflow} \rrbracket \nonumber \\
%\mathbf{Dataflow} ::=\; & \mathbf{0} \mid x \mid l_D <\bar{x}>\, \mid \mathbf{Subsys} <\bar{x}, \bar{x} > \nonumber \\
% & \mid \mathbf{Dataflow} \parallel \mathbf{Dataflow} \mid (\nu \bar{x}) \, \mathbf{Dataflow} \nonumber \\
%\mathbf{MsgType} ::=\; & M_{\bar{e},e_{ext}} \\
%\mathbf{SysTypeDef} ::=\; & \mathbf{Subsys} <\bar{i},\bar{o}> \, \mid l_S <x,x> \nonumber \\
% & \mid \mathbf{SysTypeDef} \parallel \mathbf{SysTypeDef} \mid \mathbf{MsgType} <\bar{x},y> \nonumber \\
%\mathbf{SysType} ::=\; & L_{\bar{y_i},\bar{y_o}}^{Sys} \llbracket (\nu \bar{x})(\nu \bar{e}) \, \mathbf{SysTypeDef} \rrbracket \nonumber \\
%\mathbf{LogicalModel} ::=\; & \mathbf{SysType} <\bar{i},\bar{o}>\, \mid \, l_L <o,i>  \nonumber \\
% & \mid \mathbf{LogicalModel} \parallel \mathbf{LogicalModel} \nonumber
%\end{align}
%
%Briefly (from the bottom rule to the top), logical 
%models consist of component blocks ($\mathbf{SysType}$) 
%whose interface ports connected by edges.  Component 
%blocks are specified by Simulink dataflow blocks 
%($\mathbf{Dataflow}$) whose interface ports are 
%connected either to other Simulink dataflow blocks 
%or to fields in message instances. Each message 
%instance ($\mathbf{MsgType}$) inside a system component 
%type block also has an interface vertex ($y$) which
%faces outward, and all other vertices are hidden within 
%the component $(\nu \bar{x})(\nu \bar{y}) \mathbf{SysTypeDef}$.  
%At the logical architecture model level, data is 
%exchanged via messages which aggregate the individual 
%dataflow connections within the components.  
%$\mathbf{Dataflow}$ blocks are built up from connections 
%between functional vertices and between the interfaces 
%on composite subsystem blocks ($\mathbf{Subsys}$).  
%These each correspond to sorts in the ESMoL term algebra.
%
%Let $i$, $o$, and $e$ be vertex sorts corresponding to input ports, output 
%ports, and message elements respectively.  Let $s$, $c$, $f$, and $d$ be
%edge sorts (of $l_D$, above) representing signal edges, connectivity edges
%(as described above to support the incremental interface), $f$ for Simulink
%primitive function blocks, and $d$ for delay blocks.  The $f$ function edge
%sorts are n-ary, so each function block can have an arbitrary but finite
%number of input and output connections. For $l_S$ define the
%sorts (given with their interfaces) $l^{b,b} < o, i >$, $l^{m,b} <e,i>$, 
%and $l^{b,m} <o,e>$.  These represent the three different connection types in
%a $\mathbf{SysType}$ specification, for connecting between ports of Simulink
%blocks (from outputs to inputs) ($l^{b,b}$), from message elements to Simulink 
%input ports ($l^{m,b}$), and from Simulink output ports to message elements
%($l^{b,m}$).

%Finally we give an abbreviated encoding of terms representing ESMoL models into the more general hierarchical graph algebra. The encoding assigns the various layers of hierarchy from the ESMoL component type system to hierarchical designs in the graph.  Edges from all layers map to (possibly generalized) edges in the new graph, and ESMoL ports map to vertices.
%
%\begin{align}
%%x &= x & \nonumber \\
%L_{\bar{i},\bar{o}}^{Subsys} \llbracket \mathbf{Dataflow} \rrbracket <\bar{x},\bar{x}> &= L_{\bar{x}} \llbracket \mathbb{G} \rrbracket < \bar{x} > & \nonumber \\
%%M_{\bar{e},e_{ext}} \llbracket \rrbracket <\bar{x}, y> &= L_{\bar{x}} \llbracket \mathbb{G} \rrbracket <\bar{x}> & \nonumber \\
%L_{\bar{i},\bar{o}}^{Sys} \llbracket (\nu \bar{x})(\nu \bar{e}) \mathbf{SysTypeDef} \rrbracket &=  L_{\bar{y}} \llbracket \mathbb{G} \rrbracket <\bar{x}> & \\
%l_D <\bar{x}> &= l <\bar{x}> & \nonumber \\
%l_{*} <x,x> &= l_c <\bar{x}> & \nonumber
%%(\nu \bar{x}) \mathbf{Dataflow} &= (\nu \bar{x}) \mathbb{G} & \nonumber
%\end{align}

To find delay-free loops for a given ESMoL model, we must first map a well-formed ESMoL model to the generic hierarchical graph model (as in \cite{modeling:esmol_cycles_tr}), remove all delay elements from the model, and  then invoke the algorithm.  For any cycle found in a component we can construct a more detailed cycle by substituting paths using the connectivity edges with their more detailed equivalents in the descendants of the component (recursively descending downwards until we run out of cycle elements).  Call this subgraph the \emph{expanded cycle}.  Repeating the cycle enumeration algorithm on these structures yields the full set of elementary cycles, and should still retain considerable efficiency as we are only analyzing cycles with possible subcycles, which can be a relatively small slice of the design graph.  \cite{modeling:esmol_cycles_tr} includes an example.
 

\section{Evaluation}

\subsection{Fixed-Wing Aircraft Example}

\begin{figure}[htb]	
    \begin{center}
    \centerline{\includegraphics[width=\columnwidth]{figures/fixed_wing_controller_mdl.png}}
    \caption{Simulink Fixed Wing Controller Model}
    \label{fig:fw_simulink}
    \end{center}	
\end{figure}

\begin{figure}[htb]	
    \begin{center}
    \centerline{\includegraphics[width=\columnwidth]{figures/all_in_one}}
    \caption{Synchronous data flow for Fixed Wing Controller}
    \label{fig:sdfcontroller}
    \end{center}	
\end{figure}

Our case study covers cycle analysis of the control design for a fixed-wing 
aircraft.  The Simulink model (Fig. \ref{fig:fw_simulink}) shows the four controller blocks and the sensor data handler.  The particulars of the control architecture are not important for this example, but Kottenstette covers them in detail\cite{control:fixed_wing}.   The controller has five software functions which are specified as Simulink model blocks, and a dynamics component (the Cessna plant block).  The 
\textbf{MDL2MGA} model importer creates a structural replica of the Simulink
model in the ESMoL modeling language.  We use subsystems from the replica
to specify the function of synchronous software components.  
Fig. \ref{fig:sdfcontroller} illustrates one possible configuration of the fixed wing controller components.  
In this particular configuration 
(Fig. \ref{fig:sdfcontroller}) the entire dataflow is included in one
type definition, which means that the entire system will execute together
as a single synchronous function with all blocks firing at the same rate.
This particular configuration is useful for illustration, but is not the most practical implementation choice.  

% TODO: Describe the layers when we can get it working.
%In an ESMoL logical architecture model, specified software component types are
%instantiated to form a task-level dataflow model.  
%Simulink dataflow specifications.
%Component type definitions.
%Logical architecture - instances of component types (only warnings).
%Deployment models - actual malformed model detection for local dependencies.

\subsection{Incremental Analysis Results}

Table \ref{tab:cycletimes} contains data from the analysis of
the fixed wing model.  The first pass was performed incrementally, with each
subcomponent of the top level model analyzed first.  Then the top level is
analyzed using the stored path edges in the lower models.  The table reports
two run times for the analysis of each component -- the first is the processing time required to find the abstract cycles only, and the second is the full analysis which finds the expanded cycle for each abstract cycle (enumerating possibly multiple cycles per abstract cycle). The table also displays the number of hierarchical components visited and the number of individual model elements visited, together with the number of abstract cycles found and the total number of cycles.  The table row labeled \emph{top level (incremental)} contains the results for the analysis of the top level of the model once the individual path interfaces had been created for each of its subcomponents.  
The second pass (labeled \emph{top level (full)}) analyzed the entire fixed wing model at once, reporting the same quantities.  Our assessment of the scalability of the approach is inconclusive for three reasons -- 1) the model size is moderate, so overhead is likely large enough to be a significant factor in all of the run times, 2) we would need a comparison with time taken to process a fully flattened model, including the flattening traversals, and 3) we need to find larger models for our test set.  The analyzer found 18 abstract cycles and 54 detailed cycles at the top level for both passes.  The \emph{velocity\_controller} component also contained a single abstract cycle (consisting of two detailed cycles). Note that we analyzed for all cycles rather than only delay-free cycles to assess scalability. Total runtime was roughly equivalent between the full and
incremental methods for this particular model. The results so far are promising but inconclusive as far as improved performance.

\begin{table*}[htb]
\centering
\begin{tabular}[width=\columnwidth]{ | l | l | l | l | l | l | l | }

\hline
              & \textbf{Abstract} & \textbf{Full} &                     &                          & \textbf{Abstract} & \textbf{Total} \\
\textbf{Component} & \textbf{Run} & \textbf{Run}  & \textbf{Hier.} & \textbf{Total} & \textbf{Cycles} & \textbf{Cycles} \\
              & \textbf{Time (s)} & \textbf{Time (s)} & \textbf{Comps.}  & \textbf{Elts.} & \textbf{Found} & \textbf{Found} \\
\hline \hline
alpha\_beta\_mu\_controller & 0.9 & 0.9 & 9 & 80 & 0 & 0 \\
\hline
gamma\_chi\_controller & 1.6 & 1.6 & 7 & 134 & 0 & 0 \\
\hline
gamma\_chi\_mu\_sensor & 1.3 & 1.3 & 8 & 100 & 0 & 0 \\
\hline
omega\_controller & 0.9 & 0.9 & 9 & 80 & 0 & 0 \\
\hline
velocity\_controller & 0.6 & 0.8 & 6 & 60 & 1 & 2 \\
\hline
Top level (incremental) & 2.3 & 55.1 & 1 & 21 & 18 & 54 \\
\hline
Totals & 7.6 & 60.6 & & & 19 & 872 \\
\hline
 & & & & & & \\
\hline
Top level (full) & 7.9 & 60.5 & 42 & 554 & 19 & 56 \\
\hline
\end{tabular}
\caption{Cycle analysis comparisons for the fixed wing model.}
\label{tab:cycletimes}
\end{table*}

\begin{figure}[htb]	
    \begin{center}
    \centerline{\includegraphics[width=\columnwidth]{figures/velocity_cycle.png}}
    \caption{Detail of the components involved in the cycle found in the velocity controller.}
    \label{fig:velcyclemod}
    \end{center}	
\end{figure}

\begin{figure}[htb]	
    \begin{center}
    \centerline{\includegraphics[width=\columnwidth]{figures/velocity_controller_cycle1.png}}
    \caption{Full cycle for the velocity controller.}
    \label{fig:velcycle}
    \end{center}	
\end{figure}

Figs. \ref{fig:velcyclemod} and \ref{fig:velcycle} display a subset of the \emph{velocity controller} component which contains a cycle, along with the expanded cycle for the component, in order to illustrate the cycle refinement in greater detail.
The abstract cycle search discovered the presence of a cycle within the component, but part of the cycle lies within a subcomponent (\emph{anti\_windup\_control}).  The cycle detection for \emph{anti\_windup\_control} created a single path
edge in the interface between the \emph{In1} port and the \emph{Out2} port, which corresponds to two paths within  \emph{anti\_windup\_control}.  The full 
cycle as shown (Fig. \ref{fig:velcycle}) is constructed in the analyzer, and then one more pass of Johnson's algorithm resolves the two cycles within the full cycle graph as reported in Tab. \ref{tab:cycletimes}.  The extracted cycle graph is much smaller (13 elements) than the corresponding fully flattened \emph{velocity controller} model, which would contain 60 elements.




\section{Conclusion}

%The current implementation is integrated into the ESMoL tool suite for the Generic Modeling Environment\cite{mic:gme}, but thorough scalability testing requires larger models.  See Porter et al\cite{modeling:esmol_cycles_tr} for an extended account, including a more formal presentation of the method and experimental results.

One interesting observation is the generality of the approach.  Algorithm 
\ref{alg:hcycles} very nearly captures a generic procedure for bottom-up incremental syntactic analysis of hierarchical graphical models.  Note that two small contributions may emerge from this observation 1) we have a structure to which we can adapt some other model analysis techniques for incremental operation, if an appropriate component interface can be found for the particular analysis in question, and 2) this approach could lead to a tool for efficiently specifying such analyses, from which we could generate software code to implement the analysis.



\section{Acknowledgements}
This work is sponsored in part by the National Science Foundation 
(grant NSF-CCF-0820088) and by the Air Force Office of Scientific Research, 
USAF (grant/contract number FA9550-06-0312).  The views and conclusions 
contained herein are those of the authors and should not be interpreted as 
necessarily representing the official policies or endorsements, either 
expressed or implied, of the Air Force Office of Scientific Research or the 
U.S. Government.


% ---- Bibliography ----
%
\bibliographystyle{splncs03}
\bibliography{incr_cycle_analysis}

%\end{multicols}

\end{document}
